\section{A compact toolkit for jump processes}
\label{sec:ctmc-pgf}

Only a small amount of general machinery is needed before turning to the models.
The essential objects are exponential clocks, continuous-time Markov chains,
generators and probability generating functions.  This section fixes the notation
and, in particular, the convention used for forward equations.

\subsection{Exponential clocks and competing events}

A non-negative random variable \(T\) has an exponential distribution with rate
\(\lambda>0\), written \(T\sim\operatorname{Exp}(\lambda)\), when
\begin{equation}
  \Prob(T>t)=\e^{-\lambda t},
  \qquad
  f_T(t)=\lambda\e^{-\lambda t},
  \qquad t\geq0.
  \label{eq:exponential-law}
\end{equation}
The mean waiting time is \(1/\lambda\).  The defining survival function also gives
the memoryless property
\begin{equation}
  \Prob(T>s+t\mid T>s)=\Prob(T>t),
  \qquad s,t\geq0.
  \label{eq:memoryless}
\end{equation}
Consequently, after a period during which no event has occurred, the residual
waiting time has the same law as it had initially.

The other fact used repeatedly is the exponential race.  Let
\(T_1,\ldots,T_m\) be independent exponential random variables with respective
rates \(a_1,\ldots,a_m\), and write \(T=\min_rT_r\).  Independence gives
\begin{equation}
  \Prob(T>t)=\prod_{r=1}^{m}\Prob(T_r>t)
  =\exp\!\left(-t\sum_{r=1}^{m}a_r\right),
\end{equation}
so that \(T\sim\operatorname{Exp}(a_0)\), where \(a_0=\sum_ra_r\).  Moreover,
\begin{equation}
  \Prob(T_r=T)=\int_0^\infty a_r\e^{-a_rt}
       \prod_{s\neq r}\e^{-a_st}\dd t
  =\frac{a_r}{a_0}.
  \label{eq:exponential-race}
\end{equation}
Thus the total rate determines the waiting time and the relative rates determine
which event occurs.  State-dependent versions of these two statements provide a
pathwise construction of the jump processes below.

\subsection{Continuous-time Markov chains and generators}

Let \(X=(X_t)_{t\geq0}\) take values in a countable state space \(E\).  It is a
time-homogeneous \ctmc{} if, for \(s,t\geq0\), its future conditional on the past
depends only on the present state and the transition probabilities depend only on
the elapsed time:
\begin{equation}
  \Prob(X_{s+t}=j\mid X_s=i,\,X_u,0\leq u<s)
  =\Prob(X_t=j\mid X_0=i)=:P_{ij}(t).
  \label{eq:markov-property}
\end{equation}
The transition matrices \(P(t)=(P_{ij}(t))\) form a semigroup,
\(P(s+t)=P(s)P(t)\), with \(P(0)=I\).

The generator \(Q=(q_{ij})\) records the short-time transition rates.  With the
convention that \(q_{ij}\) is the rate from \(i\) to \(j\),
\begin{equation}
  q_{ij}=\lim_{h\downarrow0}\frac{P_{ij}(h)}{h}\qquad(i\neq j),
  \qquad
  q_{ii}=-\sum_{j\neq i}q_{ij}.
  \label{eq:generator}
\end{equation}
Hence \(P_{ij}(h)=q_{ij}h+o(h)\) for \(j\neq i\), while
\(P_{ii}(h)=1+q_{ii}h+o(h)\).  If the chain is in state \(i\), it waits an
exponential time with rate
\(q_i=-q_{ii}=\sum_{j\neq i}q_{ij}\) and then moves to \(j\neq i\) with
probability \(q_{ij}/q_i\).  A state with \(q_i=0\) is absorbing.

For a test function \(f:E\to\mathbb R\), the generator acts as
\begin{equation}
  (\mathcal Lf)(i)=\sum_{j\neq i}q_{ij}\{f(j)-f(i)\}.
  \label{eq:generator-action}
\end{equation}
This form makes the event structure visible: each possible jump contributes its
rate multiplied by the corresponding increment of \(f\).

Write \(p_j(t)=\Prob(X_t=j)\) and regard
\(\boldsymbol p(t)=(p_j(t))_{j\in E}\) as a row vector.  The forward equations are
\begin{equation}
  \frac{\dd p_j}{\dd t}
  =\sum_{i\neq j}p_i(t)q_{ij}-p_j(t)\sum_{k\neq j}q_{jk},
  \qquad
  \boldsymbol p'(t)=\boldsymbol p(t)Q.
  \label{eq:forward-equations}
\end{equation}
The first sum is probability flow into state \(j\); the second is flow out.  For the
transition matrix itself, the backward and forward equations are respectively
\begin{equation}
  P'(t)=QP(t),
  \qquad
  P'(t)=P(t)Q.
  \label{eq:kolmogorov-matrix}
\end{equation}
Both hold for the time-homogeneous semigroup under the usual non-explosion and
domain conditions.  We will work almost entirely with the componentwise forward
equations \cref{eq:forward-equations}, where the convention is unambiguous.

\subsection{Probability generating functions}

If \(Z\) is a random variable on \(\Nzero\), its \pgf{} is
\begin{equation}
  G_Z(z)=\E[z^Z]=\sum_{n=0}^{\infty}\Prob(Z=n)z^n,
  \qquad |z|\leq1.
  \label{eq:univariate-pgf}
\end{equation}
The distribution and its first two moments can be recovered from
\begin{align}
  \Prob(Z=n)&=\frac{G_Z^{(n)}(0)}{n!},
  &\E[Z]&=G_Z'(1),
  &\Var(Z)&=G_Z''(1)+G_Z'(1)-G_Z'(1)^2.
  \label{eq:pgf-recovery}
\end{align}
In particular, \(G_Z(0)=\Prob(Z=0)\).  If \(Z_1,\ldots,Z_k\) are independent,
then the generating function of their sum is the product of their generating
functions.  This elementary factorisation is the analytical form of the branching
property used in \cref{sec:galton-watson}.

For a process \(X_t\in\Nzero\), the time-dependent generating function
\begin{equation}
  G(z,t)=\E[z^{X_t}]=\sum_{n\geq0}p_n(t)z^n
  \label{eq:process-pgf}
\end{equation}
often converts the infinite system \cref{eq:forward-equations} into one PDE.  For
example, if state \(n\) jumps to \(n+1\) at rate \(n\lambda\) and to \(n-1\) at rate
\(n\mu\), multiplication of the forward equations by \(z^n\) and summation give
\begin{equation}
  \frac{\partial G}{\partial t}
  =(\lambda z-\mu)(z-1)\frac{\partial G}{\partial z}.
  \label{eq:linear-bd-pgf-pde}
\end{equation}
The same operation in two dimensions uses
\begin{equation}
  G(x,y,t)=\E[x^{X_t}y^{Y_t}]
  =\sum_{i,j\geq0}p_{i,j}(t)x^iy^j.
  \label{eq:bivariate-pgf}
\end{equation}
Partial derivatives then encode factors of \(i\) and \(j\).  This is why linear
per-capita event rates lead to first-order PDEs, and why the method of
characteristics is the natural solution method in \cref{sec:characteristics}.

The generating-function manipulations below may first be read as operations on
formal power series.  Whenever analytic evaluation is used, the arguments are kept
inside the unit disc or an analytic continuation is stated explicitly.  This
distinction becomes important for the hypergeometric representation in
\cref{sec:abd-hypergeometric}.
